% ============================================================
\chapter*{Pr\'eface}
\addcontentsline{toc}{chapter}{Pr\'eface}
% ============================================================

La g\'eom\'etrie riemannienne constitue l'un des piliers fondamentaux des math\'ematiques
modernes, se situant au carrefour de l'analyse, de l'alg\`ebre et de la topologie.
N\'ee des travaux visionnaires de Bernhard Riemann dans sa c\'el\`ebre \emph{Habilitationsschrift}
de 1854, intitul\'ee \emph{\"Uber die Hypothesen, welche der Geometrie zu Grunde liegen},
cette discipline a connu un d\'eveloppement spectaculaire, culminant dans des r\'esultats
aussi profonds que la preuve de la conjecture de Poincar\'e par Perelman via le flot de Ricci.

\section*{Objectifs du cours}

Ce cours s'adresse aux \'etudiants de Master et de Doctorat en math\'ematiques
pures. Il vise \`a fournir une pr\'esentation rigoureuse et compl\`ete des fondements
de la g\'eom\'etrie riemannienne, depuis la d\'efinition des vari\'et\'es riemanniennes
jusqu'aux th\'eor\`emes de comparaison et une introduction au flot de Ricci.

Les objectifs principaux sont~:
\begin{itemize}
    \item Ma\^itriser le formalisme des connexions, de la d\'erivation covariante
          et du transport parall\`ele.
    \item Comprendre en profondeur la connexion de Levi-Civita et ses propri\'et\'es
          caract\'eristiques.
    \item \'Etudier les g\'eod\'esiques, l'application exponentielle et les propri\'et\'es
          de minimalit\'e.
    \item D\'evelopper le calcul tensoriel de la courbure~: tenseur de Riemann,
          courbure sectionnelle, courbure de Ricci et courbure scalaire.
    \item \'Etablir les grands th\'eor\`emes globaux~: Bonnet--Myers, Hadamard--Cartan,
          th\'eor\`emes de comparaison de Rauch et Toponogov.
    \item Explorer la g\'eom\'etrie des espaces sym\'etriques et des sous-vari\'et\'es.
    \item Offrir une introduction au flot de Ricci, motiv\'ee par la conjecture
          de Poincar\'e.
\end{itemize}

\section*{Pr\'erequis}

Le lecteur est suppos\'e familier avec~:
\begin{itemize}
    \item La th\'eorie des vari\'et\'es diff\'erentielles~: atlas, cartes, diff\'eomorphismes,
          fibr\'e tangent, formes diff\'erentielles.
    \item L'alg\`ebre multilin\'eaire et le calcul tensoriel de base~: produits tensoriels,
          alg\`ebre ext\'erieure, contraction.
    \item Les \'el\'ements de topologie g\'en\'erale~: compacit\'e, connexit\'e, rev\^etements.
    \item Les \'equations diff\'erentielles ordinaires~: th\'eor\`eme de Cauchy--Lipschitz,
          d\'ependance r\'eguli\`ere par rapport aux conditions initiales.
    \item Les bases de la th\'eorie des groupes de Lie et de leurs alg\`ebres de Lie.
\end{itemize}

Une connaissance pr\'ealable de la g\'eom\'etrie diff\'erentielle des courbes et surfaces
dans $\R^3$ est utile mais non indispensable~; les concepts seront revisit\'es dans le
cadre intrins\`eque.

\section*{Organisation du cours}

Le cours est organis\'e en onze chapitres, structur\'es comme suit~:

\begin{description}
    \item[Chapitre 1 --- Vari\'et\'es riemanniennes.]
    On introduit les m\'etriques riemanniennes, la structure d'espace m\'etrique
    associ\'ee, et les premiers exemples fondamentaux~: sph\`eres $S^n$, espaces
    hyperboliques $\mathbb{H}^n$, tores plats, groupes de Lie munis de m\'etriques
    bi-invariantes. On \'etudie les isom\'etries et les notions de compl\'etude.

    \item[Chapitre 2 --- Connexions et d\'eriv\'ees covariantes.]
    On d\'eveloppe la th\'eorie g\'en\'erale des connexions sur les fibr\'es vectoriels~:
    d\'efinition axiomatique, transport parall\`ele, torsion, holonomie.

    \item[Chapitre 3 --- Connexion de Levi-Civita.]
    On d\'emontre le th\'eor\`eme fondamental de la g\'eom\'etrie riemannienne~:
    existence et unicit\'e de la connexion sans torsion compatible avec la m\'etrique.
    Les symboles de Christoffel sont calcul\'es explicitement sur les exemples.

    \item[Chapitre 4 --- G\'eod\'esiques et application exponentielle.]
    On \'etudie les g\'eod\'esiques comme courbes localement minimisantes,
    l'application exponentielle, le lemme de Gauss, les coordonn\'ees normales,
    et le rayon d'injectivit\'e.

    \item[Chapitre 5 --- Courbure~: tenseur de Riemann.]
    On introduit le tenseur de courbure de Riemann via la non-commutativit\'e
    de la d\'erivation covariante, ses sym\'etries fondamentales, et les identit\'es
    de Bianchi.

    \item[Chapitre 6 --- Courbures sectionnelle, de Ricci et scalaire.]
    On d\'efinit les diff\'erentes contractions du tenseur de Riemann et on les
    calcule sur les espaces mod\`eles.

    \item[Chapitre 7 --- Th\'eor\`emes de comparaison.]
    On pr\'esente les th\'eor\`emes de Rauch, Toponogov et Bishop--Gromov,
    qui relient les bornes de courbure \`a la g\'eom\'etrie globale.

    \item[Chapitre 8 --- Espaces sym\'etriques.]
    On \'etudie les espaces sym\'etriques riemanniens au sens de Cartan,
    leur lien avec les groupes de Lie, et la classification en types compacts
    et non compacts.

    \item[Chapitre 9 --- Bonnet--Myers et Hadamard--Cartan.]
    On d\'emontre ces deux th\'eor\`emes fondamentaux qui illustrent l'influence
    du signe de la courbure sur la topologie globale.

    \item[Chapitre 10 --- G\'eom\'etrie des sous-vari\'et\'es.]
    On d\'eveloppe la th\'eorie des sous-vari\'et\'es~: seconde forme fondamentale,
    \'equations de Gauss--Codazzi--Ricci, hypersurfaces, courbure moyenne.

    \item[Chapitre 11 --- Introduction au flot de Ricci.]
    On pr\'esente le flot de Ricci de Hamilton, l'existence locale, les solutions
    autosimilaires (solitons), et un aper\c{c}u de la strat\'egie de Perelman.
\end{description}

\section*{Conventions et notations}

Tout au long de ce cours, nous adoptons les conventions suivantes~:
\begin{itemize}
    \item Les vari\'et\'es sont lisses ($C^\infty$), de Hausdorff, \`a base
          d\'enombrable, et connexes sauf mention contraire.
    \item $(M, g)$ d\'esigne une vari\'et\'e riemannienne, $\nabla$ la connexion de
          Levi-Civita associ\'ee, $R$ le tenseur de courbure.
    \item Convention de signe pour le tenseur de Riemann~:
          \[
              R(X,Y)Z = \nabla_X \nabla_Y Z - \nabla_Y \nabla_X Z - \nabla_{[X,Y]} Z.
          \]
    \item La convention de sommation d'Einstein est employ\'ee~: un indice r\'ep\'et\'e
          en positions contravariante et covariante est implicitement somm\'e.
    \item $\mathfrak{X}(M)$ d\'esigne l'espace des champs de vecteurs lisses sur $M$.
    \item $\Gamma(E)$ d\'esigne l'espace des sections lisses d'un fibr\'e vectoriel $E$.
    \item $T_pM$ d\'esigne l'espace tangent en $p$, $T^*_pM$ l'espace cotangent.
    \item $\inner{\cdot}{\cdot}$ ou $g(\cdot, \cdot)$ d\'esigne le produit scalaire riemannien.
    \item $\norm{v} = \sqrt{\inner{v}{v}}$ d\'esigne la norme riemannienne.
    \item $d(p,q)$ d\'esigne la distance riemannienne entre $p$ et $q$.
    \item $B(p,r) = \{q \in M : d(p,q) < r\}$ d\'esigne la boule g\'eod\'esique.
    \item $\operatorname{Ric}$ d\'esigne le tenseur de Ricci, $S$ ou $\operatorname{Scal}$
          la courbure scalaire.
    \item $K(\sigma)$ d\'esigne la courbure sectionnelle du $2$-plan $\sigma$.
\end{itemize}

\begin{formules}[title=\textbf{Espaces mod\`eles de r\'ef\'erence}]
Les trois familles d'espaces mod\`eles \`a courbure constante~:
\begin{enumerate}
    \item \textbf{Sph\`ere} $S^n(r)$~: courbure sectionnelle $K = 1/r^2$,
          $\pi_1 = 0$ pour $n \geq 2$.
    \item \textbf{Espace euclidien} $\R^n$~: courbure $K = 0$, plat.
    \item \textbf{Espace hyperbolique} $\mathbb{H}^n(r)$~: courbure $K = -1/r^2$,
          simplement connexe, non compact.
\end{enumerate}
\end{formules}

\section*{Fil conducteur~: courbure et topologie}

Le th\`eme central de ce cours est la relation entre la \emph{courbure} (donn\'ee
locale, analytique) et la \emph{topologie} (donn\'ee globale, qualitative). Cette
dialectique se manifeste dans les r\'esultats suivants~:

\begin{center}
\renewcommand{\arraystretch}{1.3}
\begin{tabular}{lll}
\hline
\textbf{Condition sur la courbure} & \textbf{Cons\'equence topologique} & \textbf{Th\'eor\`eme} \\
\hline
$\operatorname{Ric} \geq (n-1)\kappa > 0$ & $M$ compacte, $\pi_1$ fini & Bonnet--Myers \\
$K \leq 0$ partout & $\widetilde{M} \cong \R^n$ & Hadamard--Cartan \\
$K \geq \delta > 0$ & $\operatorname{diam}(M) \leq \pi/\sqrt{\delta}$ & Bonnet--Myers \\
$K > 1/4$ pincement & $M$ hom\'eomorphe \`a $S^n$ & Sph\`ere diff.\\
\hline
\end{tabular}
\end{center}

\section*{Conseils de travail}

La g\'eom\'etrie riemannienne est une discipline o\`u l'intuition g\'eom\'etrique et la
rigueur analytique doivent se compl\'eter. Nous recommandons~:
\begin{itemize}
    \item De toujours v\'erifier les formules g\'en\'erales sur les exemples concrets
          ($S^2$, $\mathbb{H}^2$, tore plat).
    \item De dessiner les situations g\'eom\'etriques, m\^eme schématiquement.
    \item De r\'esoudre les exercices propos\'es, qui font partie int\'egrante du cours.
    \item De comparer syst\'ematiquement les approches intrins\`eque et en coordonn\'ees.
\end{itemize}

\section*{R\'ef\'erences bibliographiques}

\begin{enumerate}
    \item M.~P.~do Carmo, \emph{Riemannian Geometry}, Birkh\"auser, 1992.
    \item J.~M.~Lee, \emph{Riemannian Manifolds: An Introduction to Curvature},
          Springer, 1997.
    \item S.~Gallot, D.~Hulin, J.~Lafontaine, \emph{Riemannian Geometry},
          Universitext, Springer, 2004.
    \item P.~Petersen, \emph{Riemannian Geometry}, GTM 171, Springer, 2016.
    \item J.~Jost, \emph{Riemannian Geometry and Geometric Analysis}, Springer, 2017.
    \item J.~Cheeger, D.~Ebin, \emph{Comparison Theorems in Riemannian Geometry},
          AMS Chelsea, 2008.
    \item S.~Kobayashi, K.~Nomizu, \emph{Foundations of Differential Geometry},
          vol.~I et II, Wiley, 1963--1969.
    \item B.~Chow, P.~Lu, L.~Ni, \emph{Hamilton's Ricci Flow}, AMS, 2006.
    \item M.~Berger, \emph{A Panoramic View of Riemannian Geometry}, Springer, 2003.
\end{enumerate}

\bigskip
\hfill\emph{L'auteur, mars 2026}
